% Generated draft: requires mathematical and source review.
\documentclass[11pt]{article}
\usepackage[margin=25mm]{geometry}
\usepackage{fontspec}
\setmainfont{Latin Modern Roman}
\usepackage{amsmath,amssymb,mathtools,unicode-math}
\setmathfont{Latin Modern Math}
\usepackage{enumitem,needspace,xurl,hyperref}
\setlength{\parindent}{0pt}
\setlength{\parskip}{5pt}
\setlength{\emergencystretch}{3em}
\newcommand{\N}{\mathbb{N}}
\newcommand{\Z}{\mathbb{Z}}
\newcommand{\Q}{\mathbb{Q}}
\newcommand{\R}{\mathbb{R}}
\newcommand{\C}{\mathbb{C}}
\newcommand{\F}{\mathbb{F}}
\newcommand{\A}{\mathbb{A}}
\newcommand{\PP}{\mathbb{P}}
\newcommand{\E}{\mathbb{E}}
\newcommand{\eps}{\varepsilon}
\newcommand{\dd}{\,\mathrm{d}}
\begin{document}
\section*{UnsolvedMath ID: 20003087}
\textbf{Source problem code:} catalogue entry
\subsection{Definitions and mathematical statement}
Let \(S=S_{g,p}\) be an orientable finite-type surface with \(2g-2+p>0\), and let \(X\in\mathcal T(S)\) be a complete finite-area marked hyperbolic structure. For a tangent vector \(v\), the asymmetric Thurston infinitesimal norm is
\[
\|v\|_{\rm Th,X}=\sup_{\gamma}\frac{\dd\ell_\gamma|_X(v)}{\ell_\gamma(X)},
\]
where \(\gamma\) ranges over essential nonperipheral simple closed curve classes, or equivalently the appropriate nonzero measured laminations. The pair-of-pants case has zero-dimensional Teichm\"uller space and is handled separately. An abstract normed tangent space is compared by a real linear bijection preserving this asymmetric norm, without replacing it by a symmetrization. The question attributed to K. Rafi asks which properties of \(X\) this abstract linear-isometry class determines.

The supplied aid attributes to Pan's 2023 and Huang--Ohshika--Papadopoulos's 2025 rigidity theorems that, in the stated complete finite-area finite-type setting, it determines the topological type and the full unoriented unmarked hyperbolic isometry type. All low-complexity ranges and ineffective mapping-class kernels must be checked in those theorems; the record does not reproduce them. In particular it does not determine a preferred marking, orientation or preassigned labels for curves.

The aid packages the recovered data as the extended-mapping-class orbit of the length-weighted curve complex: essential simple closed curves carry weights \(\ell_\gamma(X)\), with their disjointness and topological-type data, and the entire labeling is considered only up to the extended mapping class action. This is more precise than arbitrary relabeling of a bare weighted set, especially in low complexity. It reports equivalence of information between this orbit, the unoriented isometry type and the asymmetric tangent norm.

Consequently the reported recoverable invariants include \((g,p)\), area \(2\pi(2g-2+p)\), the unmarked simple length spectrum, the systole and the incidence complex of systolic curves, curve-type-refined systolic lengths, pants-decomposition length data, and the surface's full isometry group. For example the shortest-pants maximum-length invariant can be written
\[
\inf_{\mathcal P\text{ a pants decomposition}}\max_{\gamma\in\mathcal P}\ell_\gamma(X),
\]
when pants curves are present, with a separate convention for a pair of pants. Finally a norm-automorphism group is not automatically the entire surface isometry group: under the reported rigidity correspondence it is the image of that group on the tangent space, hence its quotient by the kernel of the derivative action. The aid mentions exact low-complexity kernel corrections but does not list them; no unsupported identification of a faithful action is made. These invariant extractions are consequences of the supplied rigidity account, not independent proofs of those theorems.
\subsection{Short English statement}
Determine the geometric information in the asymmetric Thurston tangent norm. The supplied rigidity account says it recovers the unmarked unoriented hyperbolic surface, hence its length and systolic data, while leaving markings and orientation undetermined and requiring low-complexity kernel corrections for automorphism groups.
\subsection{Sources}
\begin{itemize}
\item K. Rafi, Problem 1.2, AIM workshop \emph{Lipschitz metric on Teichmueller space}, \url{https://aimath.org/pastworkshops/lipschitzteichproblems.pdf}. The word-break artifact in ``determined'' and decorative dots are omitted without changing the question.
\item The research summary cites Pan (2023) and Huang--Ohshika--Papadopoulos (2025) for infinitesimal rigidity, without full publication titles, URLs, theorem statements or exceptional-kernel lists.
\item Local UnsolvedMath registry, permanent ID 20003087, exact code \texttt{catalogue entry}; canonical location \texttt{aim-topology-notes.json}, notes[174], dated 14 August 2026; statement labeled exact, registry status solved and research classification solved in literature.
\item Research attempt 1: full solution, problem solved at the run, rigor label proved, novelty confidence low. The weighted-curve-complex and invariant ledger are machine-generated research aids, not peer reviewed; independent verification is required before citation.
\end{itemize}
\end{document}
